\bookchapter{פרק 45: הבטחת איכות — לגרום לקוד שה-AI כותב להיות אמין}

\chapterintro{לקוד שנכתב על ידי AI יש מאפיין אחד: הוא בדרך כלל “נראה נכון” — שמות משתנים תקניים, הערות מלאות, לוגיקה ברורה. אבל “נראה נכון” לא שווה ל“באמת נכון”. הבעיות החבויות בקוד של AI קשות לגילוי יותר מאשר בקוד שנכתב ביד, כי איכות החזית שלו גבוהה מדי וגורמת לנו להירגע. בפרק זה תלמדו כיצד לבנות לקוד AI מערכת הבטחת איכות אפקטיבית — שלוש חזיתות הגנה של בדיקות, סקירה ובקרת גרסאות — כדי להפוך את “נראה נכון” ל“באמת נכון”.}

\booksection{א. תכנות מונחה-בדיקות: תנו ל-AI לכתוב את הבדיקות}

הבדיקות הן חזית ההגנה הראשונה להבטחת איכות בתכנות עם AI. החדשות הטובות: לבקש מ-AI לכתוב בדיקות קל ואמין יותר מאשר לבקש ממנו לכתוב קוד פונקציונלי — כי ההתנהגות הצפויה של בדיקה מוגדרת היטב, ופחות סיכוי לקבל קוד ש“נראה נכון אבל למעשה אינו נכון”.

תהליך התכנות המונחה-בדיקות עם AI הוא כזה: קודם תנו ל-AI לכתוב את קוד הפונקציונליות; מרגע שאימתתם שהתפקוד הבסיסי עובד, בקשו ממנו מיד לכתוב מקרי בדיקה. מקרי הבדיקה צריכים לכסות את המסלול התקין, מקרי קצה ומקרים חריגים. לאחר מכן תנו ל-AI לתקן את הקוד הפונקציונלי בהתאם למשוב הבדיקות. היתרון בתהליך הזה: ברגע שהבדיקות נכתבו, הן הופכות לכלי אימות מתמשך — כל שינוי עתידי אפשר לאמת על ידי הרצת הבדיקות, כדי לוודא שלא נכנסו בעיות רגרסיה.

\begin{lstlisting}
כתוב בדיקות יחידה לפונקציה calculate_discount.
מקרי הבדיקה צריכים לכסות:
1. מסלול תקין:
   - 100 יואן עם 20% הנחה -> 80
   - 200 יואן עם 50% הנחה -> 100
2. מקרי קצה:
   - 0 יואן עם 20% הנחה -> 0
   - 100 יואן בלי הנחה -> 100 (אין הנחה)
   - 100 יואן במחיר מלא -> 100 (מחיר מלא)
3. מקרים חריגים:
   - מחיר שלילי -> צריך לזרוק שגיאה
   - אחוז ההנחה > 1 -> צריך לזרוק שגיאה
   - קלט שאינו מספר -> צריך לזרוק שגיאה
השתמש בפריימוורק Vitest ותן את קובץ הבדיקות המלא.
\end{lstlisting}

שימו לב למאפיין של ההנחיה הזאת: היא מפרטת במפורש שלוש קטגוריות של מקרי בדיקה (תקין, קצה, חריג) ולא משאירה ל-AI מרחב ל“בדיקה סלקטיבית”. לפעמים AI יכתוב רק בדיקות למסלול התקין כדי “לסיים מהר”; פירוט מפורש של כל הקטגוריות מונע את הבעיה הזאת.

\booksection{ב. טכניקות לסקירת קוד}

גם אם יש כיסוי בדיקות, עדיין נדרשת סקירה אנושית של הקוד ש-AI מייצר. סקירת קוד אינה קריאה שורה-שורה — זה אורך זמן רב מדי ותלוי מדי ברמת התכנות שלכם. סקירה אפקטיבית של קוד AI צריכה להתמקד בממדי המפתח הבאים:

\begin{center}
\begin{tabularx}{\linewidth}{|X|X|X|}
\hline
\textbf{ממד הסקירה} & \textbf{על מה להקפיד} & \textbf{בעיות נפוצות} \tabularnewline
\hline
אבטחה & אימות קלט, טיפול במפתחות גישה, הגנה מפני SQL Injection & AI שוכח לעיתים קרובות לאמת קלט, ומשאיר מפתחות גישה מקודדים קשיח בתוך הקוד \tabularnewline
טיפול בשגיאות & האם חריגות נתפסות, האם הודעות השגיאה הגיוניות & AI אוהב בלוקי catch ריקים או רק console.log \tabularnewline
ביצועים & האם יש בעיות ביצועים בולטות, כמו שאילתות N+1 או לולאות אינסופיות & לפעמים AI כותב אלגוריתמים בסיבוכיות O(n²) \tabularnewline
תחזוקתיות & האם מבנה הקוד ברור, האם השמות משמעותיים & לפעמים AI מייצר קוד מופשט מדי \tabularnewline
ניהול תלויות & האם התלויות שנוספו הכרחיות, האם הגרסאות מוגדרות במפורש & AI נוטה להוסיף ספריות מיותרות \tabularnewline
\hline
\end{tabularx}
\end{center}

\begin{center}\small\color{gray}הממדים המרכזיים לסקירת קוד AI\end{center}

בסקירה אין צורך לבדוק הכול בעצמכם — אפשר לבקש מ-AI לבצע את סבב הסקירה הראשון. הנחיה לדוגמה: “בבקשה סקור את האבטחה, הטיפול בשגיאות והביצועים של הקוד הזה, וציין בעיות פוטנציאליות.” יכולת הסקירה העצמית של AI משתפרת כל העת, ולמרות שהיא אינה מחליפה סקירה אנושית באופן מלא, היא עוזרת לסנן את רוב הבעיות הבולטות.

\booksection{ג. בקרת גרסאות: התפקיד של Git בתכנות עם AI}

התפקיד של Git בתכנות עם AI חשוב עוד יותר מאשר בתכנות מסורתי. מכיוון שכל שינוי של AI עלול להכניס שינויים בלתי צפויים, אתם זקוקים ל“תרופת חרטה” אמינה — ובקרת הגרסאות היא בדיוק אותה תרופת חרטה.

שיטת העבודה המיטבית עם Git בתכנות עם AI היא “commit תדיר, התקדמות בצעדים קטנים”. בכל פעם שמושלמת נקודת פונקציונליות שניתנת לאימות — בצעו commit. כך, כשאחד השינויים של AI משתבש, אפשר לחזור במדויק למצב התקין הקודם בלי להפיל את כל העבודה.

טריק שימושי נוסף הוא “בידוד בענפים”. לשינויים של AI שתוצאתם אינה ודאית (למשל רפקטורינג גדול), פתחו קודם ענף חדש ועבדו שם. אם התוצאה משביעת רצון — מזגו בחזרה לענף הראשי; אם לא — פשוט מחקו את הענף, והענף הראשי לא נפגע. זה הרבה יותר בטוח מאשר לשנות שוב ושוב בענף הראשי.

\begin{enumerate}
  \item לפני כל שינוי של AI, ודאו שסביבת העבודה נקייה (אין שינויים שלא בוצע עליהם commit). כך, אם השינוי של AI לא מוצא חן בעיניכם, אפשר לחזור אחורה במהירות עם git checkout .
  \item אחרי כל שינוי של AI, ומיד כשהאימות עובר — בצעו commit. הודעת ה-commit תתאר מה השינוי הזה עשה, למשל "feat: הוספת פונקציית רישום משתמשים".
  \item לרפקטורינג גדול, פתחו קודם ענף: git checkout -b refactor/auth-system. השלימו את השינויים בענף, ומזגו רק אחרי שהבדיקות עוברות.
  \item פתחו הרגל לעיין ב-diff. אחרי כל שינוי של AI, השתמשו ב-git diff כדי לראות בדיוק מה שונה. זה לא רק עוזר לגלות בעיות — זה גם עוזר לכם ללמוד את הקוד.
\end{enumerate}

\booksection{ד. מלכודות נפוצות וכיצד להימנע מהן}

בתכנות עם AI יש כמה מלכודות שחוזרות על עצמן, וכמעט כל מתחיל נופל בהן. היכרות מוקדמת איתן תחסוך לכם הרבה עיקולים מיותרים בדרך.

\booksubsection{מלכודת ראשונה: מלכודת האמון}

“הקוד ש-AI כתב נראה מקצועי, בטח אין בו בעיה?” — זו המלכודת המסוכנת ביותר. הקוד ש-AI מייצר מסודר מאוד על פני השטח — שמות תקניים, הערות מלאות, מבנה ברור — ו“המראה המקצועי” הזה מרגיע את הערנות. אבל סדר חיצוני אינו מעיד על נכונות הלוגיקה. התרופה: אמתו תמיד באמצעות בדיקות, ולא לפי “איך זה נראה”.

\booksubsection{מלכודת שנייה: מלכודת הנפיחות}

“מכיוון ש-AI כותב לי את הקוד, למה שלא נעשה את הפיצ'רים מקיפים יותר?” — זה מוביל לנפיחות קוד. AI מצטיין ב“לתת יותר”: תבקשו ממנו עמוד התחברות, והוא ייתן לכם מערכת ניהול משתמשים שלמה. פיצ'רים מיותרים לא רק מגדילים את המורכבות, אלא גם מעלים את הסיכוי ל-bug-ים. התרופה: הקפידו על עקרון Simplicity First, ואמרו ל-AI במפורש “עשה רק את הפיצ'ר הזה, אל תוסיף נוספים”.

\booksubsection{מלכודת שלישית: מלכודת ההבנה}

“AI בטח מבין מה התכוונתי?” — לא בהכרח. יכולת ההבנה של AI חזקה, אבל מול דרישות מעורפלות הוא עדיין “מנחש”. כשהוא מנחש לא נכון — הקוד שמתקבל שגוי. והגרוע מכך: קוד שנוחש לא נכון בדרך כלל “נראה הגיוני”, מה שגורם לכם לחשוב בטעות שהוא הבין. התרופה: ספקו בהנחיה מספיק פרטים קונקרטיים, כדי לצמצם את מרחב ה“ניחושים” של ה-AI.

\booksubsection{מלכודת רביעית: מלכודת ההקשר}

“אמרתי את זה בשיחה קודמת, AI בטח זוכר?” — סביר להניח שלא. ככל שהשיחה מתארכת או שנפתחת שיחה חדשה, ה-AI מאבד את ההקשר המוקדם. אם אילוץ מרכזי לא נרשם בקובץ זיכרון הפרויקט, סביר שהוא יוזנח בשינויים הבאים. התרופה: רשמו את האילוצים המרכזיים בקובץ זיכרון הפרויקט, ואל תסתמכו על זיכרון השיחה.

\begin{tcbwarn}{רשימת זכירה לארבע המלכודות}
מלכודת האמון — אל תסתמכו על “נראה נכון”; מלכודת הנפיחות — AI אוהב “לתת יותר”; מלכודת ההבנה — מול דרישה מעורפלת הוא “מנחש”; מלכודת ההקשר — זיכרון השיחה “שוכח”. התרופות המתאימות, בהתאמה: אימות בבדיקות, חיתוך המיותר, מתן פרטים מלאים, ורישום בקובץ הזיכרון.

\end{tcbwarn}

\begin{tcbpractice}{היד על הדופק: הקימו שלוש חזיתות הגנה לפרויקט שלכם}
אם כבר התחלתם את פרויקט מספר 1 (אתר הכלים): (1) בקשו מ-AI לכתוב שלוש קטגוריות של מקרי בדיקה לפונקציות הליבה; (2) סקרו את הקוד לפי חמשת הממדים (אבטחה / טיפול בשגיאות / ביצועים / תחזוקתיות / ניהול תלויות); (3) אתחלו את הפרויקט עם git ובצעו את ה-commit הראשון. רק כששלוש חזיתות ההגנה שלמות — הפרויקט שלכם באמת שלכם.

\end{tcbpractice}

\begin{tcbsummary}{סיכום הפרק}
\begin{itemize}
  \item קוד AI ש"נראה נכון" $\neq$ "באמת נכון" — דווקא איכות חזית גבוהה מדי מרגיעה את הערנות
  \item תנו ל-AI לכתוב בדיקות: פרטו במפורש שלוש קטגוריות — תקין / קצה / חריג — כדי למנוע "בדיקה סלקטיבית"
  \item חמישה ממדי סקירה: אבטחה, טיפול בשגיאות, ביצועים, תחזוקתיות, ניהול תלויות
  \item Git הוא "תרופת החרטה": commit תדיר, צעדים קטנים, בידוד בענפים, ועיון קבוע ב-diff
  \item ארבע המלכודות: אמון (אימות בבדיקות), נפיחות (חיתוך המיותר), הבנה (פרטים מלאים), הקשר (רישום בקובץ זיכרון)
\end{itemize}

\end{tcbsummary}
